\chapter{Estabilidad Global: El Teorema de la Salida}
Uno de los resultados centrales de TCCU-6D es el \emph{Teorema de la Salida}, que garantiza que el universo no colapsa en una singularidad ni se fragmenta en regiones inconexas, sino que alcanza un estado estable de coherencia finita.

\section{Enunciado del teorema}
Sea $\Phi$ una solución de las ecuaciones de campo con condiciones iniciales regulares. Entonces, para todo $t > 0$, $\Phi$ permanece acotado y la energía total converge a un mínimo local asintóticamente. Además, no existen singularidades desnudas.

\section{Esbozo de la demostración}
La demostración se basa en la existencia de una función de Lyapunov:
\begin{equation}
\mathcal{L} = \frac{1}{2}\dot{\Phi}^2 + \frac{1}{2}|\nabla\Phi|^2 + V(\Phi)
\end{equation}
Cuya derivada temporal es no positiva debido al término disipativo. La combinación de disipación y memoria asegura la convergencia a un atractor de dimensión finita.

\section{Implicaciones cosmológicas}
El teorema implica que el universo es \emph{antifrágil}: pequeñas perturbaciones no lo destruyen, sino que lo llevan a estados de mayor coherencia. Esto explica la estabilidad observada del cosmos a gran escala.
